\section{Reproducibility checklist details}

The exact diff that fixes the leakage in
\texttt{apps/neural\_refine/src/model.py} is:
\begin{verbatim}
- baseline_residual = baseline_data - targets_px
+ baseline_residual = baseline_data
\end{verbatim}
plus replacing per-split z-score statistics with statistics computed on
the training split alone. After this change the original ``Multi
Baseline Neural Refinement'' code reproduces the honest numbers in
Table~\ref{tab:judo} and Table~\ref{tab:self_collected}.

\section{Full ablation tables on JuDo1000}
\label{app:ablations}

We report each axis individually. All runs are with the honest pipeline,
seed $1047$, $200$ epochs of AdamW, batch size $64$, learning rate
$3 \times 10^{-4}$.

\subsection{Selection strategy}

\begin{table}[h]
\centering
\small
\begin{tabular}{lr}
\toprule
selection & test L2 (px) \\
\midrule
all 11 baselines        & 60.55 \\
top-$M$=1               & 26.42 \\
top-$M$=2               & 27.31 \\
top-$M$=3               & 28.85 \\
top-$M$=4 (paper)       & 30.46 \\
top-$M$=6               & 36.11 \\
top-$M$=8               & 41.07 \\
oracle top-$M$=4 (test) & 31.94 \\
oracle top-$M$=2 (test) & 28.04 \\
\bottomrule
\end{tabular}
\caption{Top-$M$ selection sweep on JuDo1000. Smaller $M$ is better,
because the worst classical baselines (TPS, RBF $s=0$) inject more
noise than signal. None of the configurations beats the similarity
baseline at $21.66$~px.}
\end{table}

\subsection{Noise injection $\sigma_{\max}$}

\begin{table}[h]
\centering
\small
\begin{tabular}{lr}
\toprule
$\sigma_{\max}$ (px) & test L2 (px) \\
\midrule
0  & 31.77 \\
5  & 30.94 \\
10 & 30.61 \\
20 (paper) & 30.46 \\
40 & 30.93 \\
80 & 32.05 \\
\bottomrule
\end{tabular}
\caption{Noise injection sweep. Effect is within $\pm 1$~px of the no-noise
baseline.}
\end{table}

\subsection{Network capacity}

\begin{table}[h]
\centering
\small
\begin{tabular}{lr}
\toprule
hidden dims & test L2 (px) \\
\midrule
$[16]$              & 30.71 \\
$[64]$              & 30.55 \\
$[64, 32]$          & 30.42 \\
$[64, 32, 16]$ (paper) & 30.46 \\
$[256, 128]$        & 30.84 \\
$[1024, 512]$ (released code) & 31.30 \\
\bottomrule
\end{tabular}
\caption{Capacity sweep. All within $\pm 1$~px.}
\end{table}

\subsection{Drop-one baseline-pool}

\begin{table}[h]
\centering
\small
\begin{tabular}{lr}
\toprule
removed family & test L2 (px) \\
\midrule
none (all 11)              & 60.55 \\
similarity                  & 65.21 \\
poly                        & 64.13 \\
RBF $s\in\{0,1,2\}$         & 35.92 \\
TPS                         & 41.40 \\
sim+RBF $s\in\{0,1,2\}$     & 38.05 \\
sim+TPS                     & 43.21 \\
sim+PWA                     & 42.78 \\
\bottomrule
\end{tabular}
\caption{Drop-one ablation. Removing the bare RBF and TPS families
substantially improves the honest neural variant because those baselines
contribute high-variance, low-quality predictions.}
\end{table}

\subsection{Seed sensitivity}

\begin{table}[h]
\centering
\small
\begin{tabular}{lr}
\toprule
seed & test L2 (px) \\
\midrule
1047 & 30.46 \\
7    & 31.13 \\
42   & 30.79 \\
123  & 30.65 \\
2024 & 31.04 \\
\bottomrule
\end{tabular}
\caption{Seed sweep. Mean $30.81$, std $0.27$~px across 5 seeds.}
\end{table}

\section{Per-subject heatmap}
Figure~\ref{fig:heatmap} shows per-subject test L2 across all classical
methods plus the learned shrinkage at $K{=}12$. Cells are coloured by
$\Delta$ vs the per-subject sim+RBF anchor.

\begin{figure}[h]
\centering
\includegraphics[width=\linewidth]{figA5_heatmap.pdf}
\caption{Per-subject test L2 (px) across classical methods and our
shrinkage at $K=12$. Color encodes $\Delta$ vs the per-subject
sim+RBF anchor; greens indicate improvement.}
\label{fig:heatmap}
\end{figure}
